\documentclass[11pt,a5paper,landscape]{article}

\usepackage{main}
\usepackage{exercices}

% --------------------------------------------------
% Informations du document
% --------------------------------------------------
\renewcommand{\typedocument}{Exercices}
\renewcommand{\classe}{6e}
\renewcommand{\titre}{Exercices - Vitesse}

\begin{document}

\setlist[enumerate]{topsep=0pt,itemsep=10pt,parsep=0pt,partopsep=0pt,leftmargin=*}

\setlength{\columnsep}{1cm}

\begin{multicols*}{2}

\setcounter{numexercice}{5}

% ----- Exercice 6 -----

\exercice{}

Le tableau ci-dessous contient des valeurs manquantes.

Utilise le triangle des formules pour calculer chaque valeur manquante dans la bonne unité, puis complète le tableau.

\renewcommand{\arraystretch}{1.5}
\begin{tabular}{|>{\centering\arraybackslash}m{3cm}|>{\centering\arraybackslash}m{3cm}|>{\centering\arraybackslash}m{3cm}|}
\hline
\rowcolor{gray!30}
\textbf{Vitesse} & \textbf{Distance} & \textbf{Temps} \\
\hline
\SI{50}{\meter/\second}    & ? \SI{}{\meter}              & \SI{20}{\second} \\
\hline
\SI{10}{\meter/\second}    & \SI{20}{\meter}            & ? \SI{}{\second} \\
\hline
? \SI{}{\kilo\meter/\hour}   & \SI{3}{\kilo\meter}        & \SI{15}{\minute} \\
\hline
\SI{1.5}{\meter/\second}   & \SI{15}{\meter}            & ? \SI{}{\second} \\
\hline
\SI{60}{\kilo\meter/\hour} & ? \SI{}{\kilo\meter}         & \SI{3}{\hour} \\
\hline
? \SI{}{\meter/\second}      & \SI{1}{\kilo\meter}        & \SI{100}{\second} \\
\hline
\end{tabular}

% ----- Exercice 7 -----

\exercice{Méthode du triangle}

\question{Quelle est la formule pour calculer la vitesse ?}

\question{Quelle est la formule pour calculer la masse volumique ($\rho$) ?}

\question{Pour chacune des formules, dessiner le triangle correspondant.}

\question{Pour chacune des situations suivantes :}

\begin{itemize}
    \item Identifier les éléments du triangle présents dans l'énoncé, et les réécrire avec la bonne unité.
    \item Utiliser la méthode du triangle pour trouver la formule à utiliser.
    \item Calculer la grandeur demandée et répondre à la question.
\end{itemize}

\begin{enumerate}[label=\alph*)]
    \item Une fusée avance avec la vitesse de \SI{300}{\meter/\second}. Quelle distance a-t-elle parcourue en 1 minute ?
    \item Une sculpture en pierre a un volume de \SI{90}{\centi\meter\cubed} et une masse de \SI{216}{\gram}. Quelle est sa masse volumique ?
    \item L’arrêt de bus est à \SI{0,4}{\kilo\meter}. En marchant à \SI{2}{\meter/\second} combien de temps faut-il pour y arriver ? Le prochain bus arrive dans \SI{3}{\minute}, arriveras-tu à temps ?
    \item Une voiture roule \SI{300}{\kilo\meter} en \SI{2}{\hour}. Quelle est sa vitesse ? La limitation est à \SI{130}{\kilo\meter/\hour}, est-ce que le conducteur risque une amende ?
    \item Un avion vole à \SI{900}{\kilo\meter/\hour}. Combien de temps met-il pour parcourir \SI{4500}{\kilo\meter} ?
    \item L'huile d'olive a une masse volumique de \SI{0,8}{\gram/\centi\meter\cubed}. Quelle est la masse d'une bouteille de \SI{100}{\centi\meter^3} ?
\end{enumerate}

\end{multicols*}
\end{document}